% Modelo comentado de Trabalho Acadêmico UFC em LaTeX — v3
% Classe institucional: abntexto-ufc / abntexto 1.1+

\DocumentMetadata{
  lang = pt-BR,
  pdfstandard = A-2b,
  pdfversion = 1.7
}

\documentclass{abntexto-ufc}

\ufcsetup{
  type = undergraduate-capstone,
  print-mode = single-sided,
  cover = auto,
  catalog-card = false,
  coat-of-arms = true,
  font = times,
  strict-font = false,
  tables = tabularray,
  code = listings,
  algorithms = algpseudocodex,
  glossary = glossaries,
  index = imakeidx,
  institution = {Universidade Federal do Ceará},
  center = {Nome do Centro, Faculdade, Instituto ou Campus},
  % Optional: leave blank when the work is not tied to a department or academic unit.
  department = {},
  undergraduate-program = {Nome do Curso},
  undergraduate-degree = {bacharel},
  author = {Nome Completo do Autor},
  title = {Modelo comentado de trabalho acadêmico da UFC},
  subtitle = {Guia visual de estrutura, normalização e recursos do abntexto-ufc},
  location = {Fortaleza},
  year = {2026},
  approval-date = {dia de mês de 2026},
  advisor = {Prof. Dr. Nome do Orientador},
  advisor-institution = {Universidade Federal do Ceará (UFC)},
  advisor-unit = {Nome do Centro ou Unidade},
  advisor-feminine-label = false,
  coadvisor = {},
  coadvisor-institution = {},
  coadvisor-unit = {},
  coadvisor-feminine-label = false,
  epigraph-author = {Autor da epígrafe},
  examiner-2 = {Profa. Dra. Nome do Segundo Membro},
  examiner-2-unit = {Departamento ou Unidade Acadêmica},
  examiner-2-institution = {Universidade Federal do Ceará (UFC)},
  examiner-3 = {Prof. Dr. Nome do Terceiro Membro},
  examiner-3-unit = {Departamento ou Unidade Acadêmica},
  examiner-3-institution = {Universidade Federal do Ceará (UFC)},
  examiner-4 = {Profa. Dra. Nome do Quarto Membro},
  examiner-4-unit = {Centro, Faculdade, Instituto ou Campus},
  examiner-4-institution = {Outra Instituição de Ensino Superior (SIGLA)},
  examiner-5 = {Prof. Dr. Nome do Quinto Membro},
  examiner-5-unit = {Programa de Pós-Graduação ou Unidade Acadêmica},
  examiner-5-institution = {Outra Instituição de Ensino Superior (SIGLA)},
  examiner-6 = {Profa. Dra. Nome do Sexto Membro},
  examiner-6-unit = {Departamento ou Unidade Acadêmica},
  examiner-6-institution = {Outra Instituição de Ensino Superior (SIGLA)}
}

\ufcAddBibliographyResource{backmatter/references.bib}
\input{backmatter/glossary}

\lstset{
  breaklines=true,
  columns=fullflexible,
  keepspaces=true,
  showstringspaces=false,
  numbers=none
}

% Local documentation callouts used only by this reference document. They do
% not change the public class API or the normative contract.
\newcommand{\guianormativa}[2]{%
  \par\medskip\noindent\textbf{Base normativa — #1.} #2\par\smallskip
}
\newcommand{\guiainstitucional}[1]{%
  \par\medskip\noindent\textbf{Orientação institucional UFC.} #1\par\smallskip
}
\newcommand{\guiapolitica}[1]{%
  \par\medskip\noindent\textbf{Política do modelo.} #1\par\smallskip
}
\newcommand{\guiaexemplo}[1]{%
  \par\medskip\noindent\textbf{Exemplo.} #1\par\smallskip
}

\begin{document}

\pretextual

\ufcPrintCover
\ufcPrintTitlePage
% \ufcPrintCatalogCard{frontmatter/catalog-card}
\ufcPrintErrata{frontmatter/errata}
\ufcPrintApprovalPage
\ufcPrintDedication{frontmatter/dedication}
\ufcPrintAcknowledgments{frontmatter/acknowledgments}
\ufcPrintEpigraph{frontmatter/epigraph}
\ufcPrintSummary{frontmatter/summary}
\ufcPrintAbstract{frontmatter/abstract}

\ufcPrintListOfIllustrations
\ufcPrintListOfTables
\ufcPrintListOfCodeListings
\ufcPrintListOfAlgorithms
\ufcPrintListOfAbbreviationsAndAcronyms{frontmatter/abbreviations-and-acronyms}
\ufcPrintListOfSymbols{frontmatter/symbols}
\ufcPrintTableOfContents

\textual

\section{Introdução e uso deste modelo}\label{sec:introducao}

Este documento de referência da Universidade Federal do Ceará (UFC) é, ao mesmo tempo, um exemplo compilável e um guia visual do \gls{template} \texttt{abntexto-ufc}. Em vez de simular uma pesquisa acadêmica fictícia, as próximas seções explicam a estrutura de um trabalho acadêmico, as principais regras de apresentação e de \gls{normalizacao} exercitadas pelo template e os recursos disponíveis para produção de TCCs, monografias de especialização, dissertações, teses e projetos de pesquisa.

O arquivo principal \texttt{main.tex} concentra os metadados institucionais em \texttt{\string\ufcsetup}. As pastas \texttt{frontmatter}, \texttt{chapters} e \texttt{backmatter} separam o conteúdo por etapa documental. A classe \texttt{abntexto-ufc} centraliza a apresentação e usa \texttt{abntexto} como base editorial, preservando a separação entre conteúdo e formatação característica do \LaTeX{} \parencite{knuth,lamport1994}.

\subsection{Como interpretar as orientações deste PDF}

Ao longo do documento são usados quatro tipos de indicação explícita.

\guianormativa{normas ABNT ou referência técnica aplicável}{Identifica uma regra de apresentação cuja fonte pertence à base normativa auditada do projeto. A redação deste guia é uma paráfrase técnica; o texto integral da norma deve ser consultado na fonte autorizada.}

\guiainstitucional{Identifica uma orientação, resolução, instrução normativa ou prática institucional da Universidade Federal do Ceará que complementa a base ABNT quando aplicável ao depósito ou à apresentação do trabalho.}

\guiapolitica{Identifica uma decisão técnica ou editorial do \texttt{abntexto-ufc}. Uma política do modelo não deve ser interpretada automaticamente como exigência universal da ABNT ou da UFC.}

\guiaexemplo{Identifica conteúdo incluído apenas para demonstrar visualmente uma possibilidade do template. Exemplos podem ser removidos ou substituídos no trabalho real.}

Essa separação é importante. Por exemplo, o documento de referência declara PDF/A-2b e a suíte de release o valida com veraPDF; entretanto, o perfil PDF/A-2b específico é uma política técnica do projeto, enquanto a necessidade de PDF/A no fluxo institucional de depósito é tratada separadamente pela documentação da UFC.

\subsection{Base normativa adotada}

A apresentação geral dos trabalhos acadêmicos é organizada a partir da ABNT NBR 14724:2024, considerando a versão corrigida em 1º de abril de 2025. Citações seguem a ABNT NBR 10520:2023 e referências seguem a ABNT NBR 6023:2025 \parencite{abnt14724,abnt10520,abnt6023}.

Outras referências atuam em escopos específicos: ABNT NBR 6028:2021 para resumo e palavras-chave; ABNT NBR 6024:2012 para numeração progressiva; ABNT NBR 6027:2012 para sumário; ABNT NBR 15287:2025 para projetos de pesquisa; ABNT NBR 6034:2004 para índice; ABNT NBR 12225:2023 para lombada; e Normas de Apresentação Tabular do IBGE para tabelas numéricas \parencite{abnt6028,abnt6024,abnt6027,abnt15287,abnt6034,abnt12225,ibge1993}.

\guiainstitucional{O Sistema de Bibliotecas da UFC mantém orientações de normalização e informa que os trabalhos acadêmicos devem observar as normas vigentes aplicáveis. Guias institucionais históricos continuam úteis para requisitos próprios da UFC, mas uma referência antiga a uma edição substituída da ABNT não prevalece sobre a edição vigente adotada pelo projeto \parencite{ufc2026normalizacao,ufcguia2022}.}

\subsection{O que este documento demonstra}

O PDF foi configurado no perfil \texttt{undergraduate-capstone}, em impressão apenas no anverso, com fonte Times em modo portátil, listas, glossário, índice, códigos, algoritmos e tabelas \texttt{tabularray}. Essa configuração serve como corpus visual completo; outros perfis alteram automaticamente os elementos que lhes são aplicáveis.

A suíte de validação do repositório cobre, entre outros pontos, tamanho A4, margens, paginação, tipografia, pré-textuais, citações, referências, objetos, tabelas, equações, códigos, algoritmos, projetos, elementos pós-textuais, matriz de perfis, motores de compilação, compatibilidade com ambiente semelhante ao Overleaf e certificação de Times New Roman e Arial literais em Windows.

\guiapolitica{O template reduz o trabalho manual de formatação, mas não substitui a verificação final das exigências específicas do curso, programa, edital, repositório institucional ou procedimento administrativo aplicável ao trabalho.}

\index{normalização}\index{LaTeX}\index{ABNT}\index{UFC}

\input{chapters/2-theoretical-framework}
\section{Elementos pré-textuais em detalhe}\label{sec:pretextuais-detalhe}

Esta seção explica os principais elementos que aparecem antes do texto acadêmico. O próprio PDF já apresentou cada um deles nas páginas anteriores; aqui o objetivo é tornar explícita a finalidade, a aplicabilidade e a regra de apresentação exercitada pelo template.

\subsection{Capa e folha de rosto}

A capa identifica externamente o trabalho. A folha de rosto repete os dados essenciais e acrescenta informações acadêmicas como natureza do trabalho, curso ou programa, orientação e demais campos aplicáveis ao perfil.

\guianormativa{ABNT NBR 14724:2024}{Capa e folha de rosto integram a estrutura de trabalhos acadêmicos. A ordem e a composição dos campos são tratadas pelo perfil selecionado no \texttt{abntexto-ufc} \parencite{abnt14724}.}

\guiapolitica{Os dados pessoais e institucionais são declarados em \texttt{\string\ufcsetup}. Isso evita editar manualmente a diagramação da capa ou da folha de rosto e permite que os perfis reutilizem os mesmos metadados.}

\subsection{Ficha catalográfica}

A representação visual da ficha catalográfica é controlada pela opção \texttt{catalog-card}. No documento de referência ela permanece desativada para demonstrar o comportamento padrão adotado pelo perfil atual.

\guiainstitucional{A política de depósito da UFC vigente em 2026 tornou facultativa a representação visual da ficha catalográfica no arquivo depositado para os tipos documentais abrangidos pela Instrução Normativa Conjunta nº 2/2026/SIBI/PROGRAD/PRPPG \parencite{ufcinstrucao2_2026}. O autor deve conferir o procedimento institucional vigente no momento da entrega.}

\subsection{Errata}

A errata é usada quando, depois da produção do trabalho, é necessário registrar correções. O exemplo deste PDF apresenta a referência do trabalho seguida por uma tabela com folha, linha, texto incorreto e forma corrigida.

\guiaexemplo{Se não houver correções a registrar, remova a chamada \texttt{\string\ufcPrintErrata} do documento. O arquivo de exemplo existe para demonstrar a composição quando o elemento é necessário.}

\subsection{Folha de aprovação}

Nos perfis acadêmicos aos quais se aplica, a folha de aprovação registra autoria, título, natureza do trabalho, data de aprovação, orientação e banca examinadora. O exemplo usa uma banca extensa para exercitar a quebra e o posicionamento de múltiplos membros.

\guiainstitucional{No fluxo de depósito considerado pelo projeto, a folha de aprovação é apresentada sem reproduzir assinaturas digitalizadas. A conferência final deve considerar a orientação vigente do Sistema de Bibliotecas para TCC, dissertação ou tese \parencite{ufcdepositotcc2026,ufcdepositopos2026}.}

\subsection{Dedicatória}

A dedicatória é um elemento de natureza pessoal e não recebe título no exemplo UFC adotado pela classe.

\guiainstitucional{No escopo auditado do template, a dedicatória inicia abaixo do meio da folha, utiliza recuo esquerdo de 8 cm, fonte de 12 pt, espaçamento 1,5 e alinhamento justificado. Esses valores são exercitados por testes específicos do documento final.}

\guiaexemplo{O arquivo \texttt{frontmatter/dedication.tex} contém apenas uma frase curta para mostrar a posição do bloco. Em um trabalho real, substitua-a pela dedicatória desejada ou remova o elemento.}

\subsection{Agradecimentos}

Os agradecimentos registram pessoas e instituições que contribuíram para o trabalho. O título é apresentado em página própria, e o corpo mantém a tipografia do texto acadêmico exercitada pelo perfil.

\guiainstitucional{No escopo UFC auditado, o título AGRADECIMENTOS aparece centralizado, em letras maiúsculas, negrito e 12 pt; o corpo usa 12 pt, espaçamento 1,5 e alinhamento justificado.}

\guiainstitucional{Quando o trabalho tiver financiamento ao qual se aplique obrigação específica de agradecimento, como a Portaria CAPES nº 206/2018, o autor deve inserir o texto institucional exigido pela fonte correspondente \parencite{capes2062018}. O template não presume automaticamente a existência de financiamento.}

\subsection{Epígrafe}

A epígrafe apresenta uma citação relacionada ao trabalho. O template diferencia a apresentação de epígrafes curtas e longas no escopo institucional auditado.

\guiainstitucional{Epígrafes de até três linhas iniciam abaixo do meio da folha, usam recuo esquerdo de 8 cm, fonte de 12 pt, espaçamento 1,5 e aspas. Epígrafes com mais de três linhas usam recuo esquerdo de 4 cm, fonte de 10 pt, espaçamento simples e são apresentadas sem aspas.}

\subsection{Resumo, abstract e palavras-chave}

O resumo apresenta de forma concisa objetivo, método, resultados e conclusões relevantes do trabalho. O abstract oferece a versão em língua estrangeira prevista pelo perfil.

\guianormativa{ABNT NBR 6028:2021}{Para trabalhos acadêmicos, o contrato normativo adotado pelo projeto trabalha com resumo e abstract entre 150 e 500 palavras e exige palavras-chave ou keywords. O texto é apresentado em parágrafo único, fonte 12 pt, espaçamento 1,5, sem recuo de primeira linha e com alinhamento justificado \parencite{abnt6028}.}

\guiaexemplo{O resumo deste próprio PDF descreve o objetivo do modelo comentado. Assim, o pré-textual funciona como um exemplo semanticamente coerente, não como texto de preenchimento.}

\subsection{Listas}

O documento de referência habilita listas de ilustrações, tabelas, códigos, algoritmos, abreviaturas e siglas e símbolos para demonstrar todas as rotas suportadas. Em trabalhos reais, devem ser mantidas apenas as listas pertinentes ao conteúdo e às regras aplicáveis.

\guiapolitica{Figuras, gráficos e quadros são agregados na Lista de Ilustrações. Tabelas, códigos e algoritmos possuem listas próprias no modelo de referência. A existência dessas listas no PDF de demonstração não significa que todo trabalho precise conter todos esses tipos de objeto.}

\subsection{Sumário}

O sumário apresenta a estrutura de seções do documento e os respectivos números de página.

\guianormativa{ABNT NBR 6027:2012}{O template mantém o título SUMÁRIO centralizado e em letras maiúsculas, posiciona os números de página à direita e reproduz a hierarquia tipográfica das seções. Elementos pré-textuais não são incluídos no sumário \parencite{abnt6027}.}

\index{capa}\index{folha de rosto}\index{errata}\index{dedicatória}\index{agradecimentos}\index{epígrafe}\index{resumo}\index{sumário}

\section{Formatação geral e organização da parte textual}\label{sec:formatacao-geral}

A apresentação do texto deve ser controlada pela classe, e não por ajustes manuais espalhados pelo conteúdo. Esta seção resume as regras gerais que o documento de referência exercita e que a suíte de validação mede diretamente no PDF produzido.

\subsection{Papel e margens}

\guianormativa{ABNT NBR 14724:2024}{O documento é produzido em papel A4. No anverso, a configuração adotada utiliza margens esquerda e superior de 3 cm e margens direita e inferior de 2 cm. Em impressão frente e verso, as margens laterais são espelhadas para preservar a margem interna maior \parencite{abnt14724}.}

\guiaexemplo{Este PDF usa \texttt{print-mode=single-sided}. Os testes do repositório também geram páginas de verso e medem a geometria final para verificar o espelhamento das margens.}

\subsection{Fonte, tamanho e cor}

O corpo textual do documento usa tamanho 12. A API pública oferece as famílias \texttt{times} e \texttt{arial}; o modo \texttt{strict-font} diferencia a exigência de identidade literal da fonte e o uso de fallback portátil para desenvolvimento.

\guianormativa{ABNT NBR 14724:2024 e política tipográfica UFC adotada}{O corpo textual é apresentado em tamanho 12 e em cor preta; tamanhos reduzidos são usados em contextos específicos, como citações longas, notas, paginação e partes de objetos.}

\guiapolitica{Quando \texttt{strict-font=true}, a compilação deve falhar se a família literal solicitada não estiver disponível. Quando \texttt{strict-font=false}, o template pode usar uma família substituta explicitamente identificada para portabilidade. A certificação de Times New Roman e Arial literais é feita em uma rota Windows separada.}

\subsection{Espaçamento e parágrafos}

\guianormativa{ABNT NBR 14724:2024 e escopo UFC auditado}{O corpo usa espaçamento 1,5. Os parágrafos do corpo adotam recuo de primeira linha de 2 cm e não recebem espaço vertical adicional entre si. Contextos específicos podem usar espaço simples e tamanho reduzido.}

O recuo de primeira linha é uma propriedade do parágrafo; não deve ser simulado com espaços, tabulações ou comandos manuais no início de cada parágrafo. Da mesma forma, a separação entre parágrafos não deve ser criada com linhas em branco adicionais destinadas apenas a alterar a aparência.

\subsection{Paginação}

A contagem das folhas e a exibição dos números são controladas pela classe. No fluxo exercitado pelo template, os números tornam-se visíveis a partir da primeira página textual e usam tamanho reduzido.

\guianormativa{ABNT NBR 14724:2024}{A apresentação da paginação acompanha a progressão física do documento e respeita a distinção entre anverso e verso. O template posiciona o número no canto superior externo da página conforme o modo de impressão.}

\subsection{Seções e subseções}

\guianormativa{ABNT NBR 6024:2012}{A numeração é progressiva e a hierarquia deve ser refletida de forma consistente no corpo e no sumário \parencite{abnt6024}.}

\guiapolitica{A classe limita sua hierarquia pública exercitada a cinco níveis. Cada seção primária inicia em nova página, e a classe controla o espaçamento e a apresentação do indicador numérico.}

Evite aplicar negrito, tamanho, alinhamento ou espaçamento manualmente aos títulos para “corrigir” uma página isolada. Se um título estiver incorreto, a correção deve ocorrer na configuração ou na implementação do nível correspondente.

\subsection{Notas de rodapé}

Notas de rodapé são apresentadas com tamanho reduzido e espaço simples. A composição deve permanecer legível e claramente separada do corpo.

\guiainstitucional{No escopo auditado atual, o texto da nota usa 10 pt, espaçamento simples e alinhamento suspenso nas linhas de continuação. O separador é exercitado com 5 cm a partir da margem esquerda.}

\subsection{Alíneas e subalíneas}

Quando uma enumeração fizer parte da estrutura discursiva do texto, o template fornece ambientes para alíneas e subalíneas que preservam a apresentação adotada pela classe. Enumerações comuns do \LaTeX{} continuam disponíveis quando não representam a estrutura normativa de alíneas.

\guiaexemplo{A Seção~\ref{sec:catalogo-exemplos} apresenta alíneas, subalíneas e uma enumeração comum para comparação visual.}

\index{margens}\index{fonte}\index{espaçamento}\index{parágrafo}\index{paginação}\index{seções}\index{notas de rodapé}

\section{Citações, notas e referências}\label{sec:citacoes-referencias}

Citações conectam o argumento do autor às fontes utilizadas; referências identificam os documentos citados ou utilizados de acordo com a política bibliográfica adotada. O template separa a apresentação das citações no texto da composição final da lista de referências.

\subsection{Citação indireta}

Na citação indireta, a ideia da fonte é apresentada com redação própria. Por exemplo, a separação entre conteúdo e apresentação é uma característica central do fluxo de preparação de documentos em \LaTeX{} \parencite{lamport1994}.

\guianormativa{ABNT NBR 10520:2023}{A indicação de autoria e data deve seguir o sistema de citação adotado no documento. O projeto usa \texttt{biblatex-abnt} como infraestrutura e mantém regressões específicas para a apresentação prevista no escopo testado \parencite{abnt10520}.}

\subsection{Citação direta curta}

Citações diretas curtas permanecem integradas ao parágrafo e usam sinais de citação apropriados. Para evitar reproduzir conteúdo protegido de uma norma técnica, o exemplo abaixo é deliberadamente sintético.

\guiaexemplo{Uma citação curta pode ser demonstrada como \enquote{texto citado com poucas linhas}, seguida da indicação bibliográfica correspondente. Em um trabalho real, o conteúdo entre aspas deve reproduzir fielmente a fonte consultada e trazer os dados exigidos para localização.}

\subsection{Citação direta longa}

\guianormativa{ABNT NBR 10520:2023 e escopo UFC auditado}{Citação direta com mais de três linhas é apresentada em bloco distinto, sem aspas, com recuo esquerdo de 4 cm, fonte reduzida de 10 pt e espaçamento simples \parencite{abnt10520}.}

O documento de referência exercita fisicamente esse comportamento com um texto sintético na Seção~\ref{sec:catalogo-exemplos}. A suíte mede recuo, tamanho de fonte, espaçamento, ausência de aspas e separação do bloco em relação ao corpo. Em uma citação direta real, a indicação bibliográfica deve incluir a página ou outro localizador aplicável quando a fonte o fornecer; o texto sintético do corpus geométrico não recebe um localizador fictício.

\subsection{Citação de citação}

O recurso conhecido como citação de citação deve ser usado com parcimônia, preferindo-se a consulta à fonte original sempre que possível. Quando necessário, a apresentação de \textit{apud} segue a política bibliográfica e a NBR 10520 vigente.

\guiapolitica{O repositório possui cenários específicos para verificar a apresentação de \textit{apud}. O fato de o template suportar o recurso não elimina a responsabilidade acadêmica de consultar a fonte original quando ela estiver disponível.}

\subsection{Notas de rodapé}

Notas podem complementar uma informação sem interromper excessivamente o fluxo do parágrafo. Sua tipografia reduzida foi descrita na Seção~\ref{sec:formatacao-geral}.

\guiaexemplo{Use notas para informação realmente complementar. Referências bibliográficas não devem ser transferidas arbitrariamente para notas apenas para contornar o sistema de citações configurado no trabalho.}

\subsection{Lista de referências}

\guianormativa{ABNT NBR 6023:2025}{A elaboração das referências segue a edição vigente adotada pelo projeto. No escopo de apresentação atualmente exercitado, as entradas são alinhadas à margem esquerda, usam fonte 12 pt e espaçamento simples internamente, com separação entre referências correspondente a uma linha em branco de espaço simples \parencite{abnt6023}.}

O arquivo \texttt{backmatter/references.bib} demonstra tipos bibliográficos diferentes: livro, artigo de periódico com DOI, trabalho em evento, dissertação, relatório técnico, norma e recurso eletrônico. O objetivo é exercitar os campos mais comuns sem sugerir que todos eles devam aparecer em todo trabalho.

\guiaexemplo{Uma referência com DOI só apresenta DOI quando esse identificador pertence ao documento referenciado. De modo semelhante, URL e data de acesso são usados quando pertinentes ao tipo de fonte e à regra bibliográfica aplicável.}

\subsection{Normas técnicas como referências}

As normas citadas neste guia aparecem também nas referências para que o leitor saiba qual edição orienta cada seção. O projeto não redistribui o texto integral das normas ABNT; a comunidade UFC pode utilizar os meios institucionais autorizados para consulta das normas técnicas.

\guiainstitucional{O Sistema de Bibliotecas da UFC disponibiliza orientação sobre acesso à coleção de normas técnicas. Este PDF identifica a norma e a edição adotadas, mas não substitui a consulta à publicação oficial \parencite{ufcnormastecnicas}.}

\index{citação direta}\index{citação indireta}\index{referências}\index{apud}\index{DOI}

\section{Ilustrações, tabelas e outros objetos acadêmicos}\label{sec:catalogo-exemplos}

Esta seção funciona como catálogo visual e corpus de regressão. Cada grupo apresenta primeiro a regra ou política exercitada e, em seguida, exemplos compiláveis. Os objetos devem ser citados e discutidos no texto do trabalho real; não devem ser incluídos apenas para preencher o documento.

\subsection{Regra geral para ilustrações}

\guianormativa{ABNT NBR 14724:2024 e NBR 10520:2023}{Ilustrações devem ser identificadas de modo consistente e apresentar indicação de fonte. Quando a fonte for externa, sua indicação deve ser compatível com o sistema de citações adotado \parencite{abnt14724,abnt10520}.}

\guiapolitica{A classe vincula título, fonte e nota à largura física do objeto. Os testes verificam que esses elementos não ultrapassem horizontalmente a largura da ilustração. A tipografia do título superior e dos blocos inferiores de fonte, legenda ou nota é validada separadamente; não ajuste seus tamanhos manualmente no conteúdo.}

\subsection{Figuras em diferentes larguras e formatos}

\legend{figure}{Figura estreita com legenda curta}
\ufcSource{Elaboração própria.}
\label{fig:estreita}
\begin{ufcobject}[here]
  \centering
  \fbox{\parbox[c][2.6cm][c]{5.5cm}{\centering Figura sintética\\5,5 cm}}
\end{ufcobject}

\guiaexemplo{A Figura~\ref{fig:estreita} demonstra que uma figura estreita não deve fazer seu título ou sua fonte ocuparem arbitrariamente toda a mancha gráfica.}

\legend{figure}{Figura de largura intermediária com uma legenda deliberadamente longa para validar quebra de linha, alinhamento e limite horizontal do título do objeto}
\ufcSource{Elaboração própria.}
\ufcNote{Exemplo sintético usado para verificar título, fonte e nota em um objeto de largura intermediária.}
\label{fig:media}
\begin{ufcobject}[here]
  \centering
  \fbox{\parbox[c][3cm][c]{9cm}{\centering Figura sintética de largura intermediária\\9 cm}}
\end{ufcobject}

\legend{figure}{Figura larga próxima à largura útil da mancha gráfica}
\ufcSource{Elaboração própria.}
\label{fig:larga}
\begin{ufcobject}[here]
  \centering
  \fbox{\parbox[c][3cm][c]{0.78\linewidth}{\centering Figura sintética larga\\78\% da largura disponível}}
\end{ufcobject}

\legend{figure}{Fluxo de processamento em arquivo PNG raster}
\ufcSource{Elaboração própria.}
\ufcNote{O arquivo é incluído diretamente com \texttt{\string\includegraphics}.}
\label{fig:png-fluxo}
\begin{ufcobject}[here]
  \centering
  \includegraphics[width=0.65\linewidth]{figures/example-flow.png}
\end{ufcobject}

As Figuras~\ref{fig:media}, \ref{fig:larga} e~\ref{fig:png-fluxo} exercitam larguras progressivamente maiores, quebra de título, fonte, nota e inclusão de imagem raster.

\subsection{Imagens institucionais reais da UFC}

\legend{figure}{Campus do Pici, onde se localiza o Departamento de Computação da UFC}
\ufcSource{Universidadefederalceara (2018), Wikimedia Commons, CC BY-SA 4.0.}
\ufcNote{Vista da Lagoa do Pici no Campus do Pici.}
\label{fig:campus-pici-computacao}
\begin{ufcobject}[here]
  \centering
  \IfFileExists{figures/ufc-campus-pici.jpg}
    {\includegraphics[width=0.72\linewidth]{figures/ufc-campus-pici.jpg}}
    {\fbox{\parbox[c][4cm][c]{0.72\linewidth}{\centering Execute \texttt{make reference-assets} para baixar a fotografia licenciada do Campus do Pici.}}}
\end{ufcobject}

\legend{figure}{Reitoria da Universidade Federal do Ceará}
\ufcSource{Joelkaula (2018), Wikimedia Commons, CC BY-SA 4.0.}
\ufcNote{Fotografia externa usada sem modificação e distribuída sob licença própria; ver \texttt{figures/LICENSES.md}.}
\label{fig:reitoria-ufc}
\begin{ufcobject}[here]
  \centering
  \IfFileExists{figures/ufc-reitoria.jpg}
    {\includegraphics[width=0.62\linewidth]{figures/ufc-reitoria.jpg}}
    {\fbox{\parbox[c][5cm][c]{0.62\linewidth}{\centering Execute \texttt{make reference-assets} para baixar a fotografia licenciada da Reitoria da UFC.}}}
\end{ufcobject}

\guiapolitica{Fotografias externas de demonstração não são incorporadas aos bundles públicos/CTAN sem controle de licença. O comando \texttt{make reference-assets} baixa os arquivos de referência e verifica sua integridade; as licenças ficam registradas separadamente. Quando a fonte externa possuir paginação ou outro localizador estável aplicável, inclua também esse localizador na indicação de fonte do objeto.}

\subsection{Gráficos}

Gráficos são tratados como ilustrações e participam da Lista de Ilustrações no modelo de referência.

\legend{grafico}{Distribuição sintética de três categorias}
\ufcSource{Elaboração própria.}
\ufcNote{Valores meramente ilustrativos.}
\label{graf:barras}
\begin{ufcobject}[here]
  \centering
  \begingroup
  \setlength{\unitlength}{1pt}
  \begin{picture}(300,180)
    \put(45,25){\line(1,0){235}}
    \put(45,25){\line(0,1){130}}
    \put(40,25){\makebox(0,0)[r]{0}}
    \put(40,58){\makebox(0,0)[r]{5}}
    \put(40,90){\makebox(0,0)[r]{10}}
    \put(40,123){\makebox(0,0)[r]{15}}
    \put(40,155){\makebox(0,0)[r]{20}}
    \put(45,58){\line(1,0){235}}
    \put(45,90){\line(1,0){235}}
    \put(45,123){\line(1,0){235}}
    \put(45,155){\line(1,0){235}}
    \put(76,25){\color{black!55}\rule{35pt}{78pt}}
    \put(148,25){\color{black!55}\rule{35pt}{117pt}}
    \put(220,25){\color{black!55}\rule{35pt}{59pt}}
    \put(94,109){\makebox(0,0){12}}
    \put(166,148){\makebox(0,0){18}}
    \put(238,90){\makebox(0,0){9}}
    \put(94,12){\makebox(0,0){Categoria A}}
    \put(166,12){\makebox(0,0){Categoria B}}
    \put(238,12){\makebox(0,0){Categoria C}}
    \put(165,0){\makebox(0,0){Categoria}}
    \put(8,90){\rotatebox{90}{Valor}}
  \end{picture}
  \endgroup
\end{ufcobject}

\guiaexemplo{O Gráfico~\ref{graf:barras} é desenhado diretamente em \LaTeX{} para que a regressão visual não dependa de um arquivo externo. Em uma pesquisa real, escolha o tipo de gráfico de acordo com os dados e informe claramente escalas, unidades, categorias e fonte.}

\subsection{Quadros}

O projeto distingue quadros de tabelas numéricas. Quadros são úteis para organizar conteúdo predominantemente textual.

\legend{quadro}{Comparação de configurações editoriais do exemplo}
\ufcSource{Elaboração própria.}
\label{qua:configuracoes}
\begin{ufcobject}[here]
  \centering
  \begin{tabular}{lll}
    \hline
    Elemento & Variação A & Variação B \\
    \hline
    Figura & estreita & larga \\
    Código & sem linhas & com linhas \\
    Algoritmo & sem linhas & com linhas \\
    \hline
  \end{tabular}
\end{ufcobject}

\guiaexemplo{O Quadro~\ref{qua:configuracoes} demonstra informação textual organizada em linhas e colunas. Não use a denominação ``tabela'' apenas porque um conteúdo foi produzido com o ambiente \texttt{tabular}; a escolha do tipo deve considerar a natureza da informação apresentada.}

\subsection{Tabelas numéricas}

\guianormativa{Normas de Apresentação Tabular do IBGE, 3. ed.}{Para tabelas numéricas, o documento de referência exercita laterais abertas, ausência de grade no corpo e regras horizontais superior, de separação do cabeçalho e inferior \parencite{ibge1993}.}

\begin{table}[htbp]
\centering
\begin{tallabnttblr}
[
  caption={Indicadores sintéticos com linhas alternadas},
  label={tab:zebra},
  remark{Fonte}={Elaboração própria.},
  remark{Nota}={A alternância de fundo é opcional e usada aqui apenas como exemplo editorial.},
]
{
  colspec={X[l] X[c] X[r]},
  row{even}={bg=black!5},
}
\toprule
Indicador & Ano & Valor \\
\midrule
Alfa & 2024 & 10 \\
Beta & 2025 & 12 \\
Gama & 2026 & 15 \\
Delta & 2026 & 18 \\
\bottomrule
\end{tallabnttblr}
\end{table}

\guiapolitica{A alternância de fundo da Tabela~\ref{tab:zebra} é uma extensão editorial opcional. Ela não é tratada como regra automática da ABNT ou do IBGE. O exemplo existe para mostrar que o módulo pode aplicar estilos sem alterar as regras estruturais da tabela.}

\subsection{Código-fonte}

Código-fonte não é convertido artificialmente em uma regra ABNT específica. A classe o trata como objeto acadêmico com identificação, fonte, numeração e largura controladas, oferecendo módulos de apresentação configuráveis.

\guiapolitica{O documento de referência usa \texttt{listings}. A política de numeração das linhas pode variar conforme a finalidade do exemplo; o importante é manter legibilidade, identificação consistente e fonte do conteúdo quando aplicável.}

\lstset{language=Python,numbers=left,stepnumber=1,numbersep=8pt}
\legend{code}{Função de média em Python com números de linha}
\ufcSource{Elaboração própria.}
\label{cod:python-numerado}
\begin{ufclisting}[here]
def media(valores):
    return sum(valores) / len(valores)
\end{ufclisting}

\lstset{language=C++,numbers=none}
\legend{code}{Função de máximo em C++ sem números de linha}
\ufcSource{Elaboração própria.}
\label{cod:cpp-sem-numero}
\begin{ufclisting}[here]
int maximo(int a, int b) {
    return a > b ? a : b;
}
\end{ufclisting}

\legend{code}{Arquivo Python externo com números de linha}
\ufcSource{Elaboração própria.}
\label{cod:python-externo}
\begin{ufcobject}[here]
  \ufcInputListing[language=Python,numbers=left,numbersep=8pt]{figures/example.py}
\end{ufcobject}

\lstset{language=Java,numbers=left,stepnumber=2,numbersep=8pt}
\legend{code}{Método Java com numeração a cada duas linhas}
\ufcSource{Elaboração própria.}
\label{cod:java}
\begin{ufclisting}[here]
static int dobro(int valor) {
    int resultado = 2 * valor;
    return resultado;
}
\end{ufclisting}
\lstset{numbers=none,stepnumber=1}

Os Códigos~\ref{cod:python-numerado}, \ref{cod:cpp-sem-numero}, \ref{cod:python-externo} e~\ref{cod:java} demonstram linguagens e políticas de numeração distintas usando o mesmo módulo.

\subsection{Algoritmos}

Algoritmos em pseudocódigo são apresentados como objetos próprios da classe. O módulo recomendado no documento de referência é \texttt{algpseudocodex}.

\guiapolitica{Numeração de linhas em pseudocódigo é uma opção editorial do exemplo e não uma exigência normativa geral. A identificação e a fonte continuam sendo tratadas de forma consistente com os demais objetos do template.}

\legend{algorithm}{Máximo divisor comum com números de linha}
\ufcSource{Elaboração própria.}
\label{alg:mdc-numerado}
\begin{ufcalgorithm}[here][1]
  \While{$b \neq 0$}
    \State $r \gets a \bmod b$
    \State $a \gets b$
    \State $b \gets r$
  \EndWhile
  \State \Return $a$
\end{ufcalgorithm}

\legend{algorithm}{Seleção do maior valor sem números de linha}
\ufcSource{Elaboração própria.}
\label{alg:maior-sem-numero}
\begin{ufcalgorithm}[here][0]
  \If{$a > b$}
    \State \Return $a$
  \Else
    \State \Return $b$
  \EndIf
\end{ufcalgorithm}

\subsection{Equações}

A Equação~\ref{eq:media-exemplo} exemplifica uma expressão numerada e referenciada no texto.
\begin{equation}
  \bar{x} = \frac{1}{n}\sum_{i=1}^{n}x_i
  \label{eq:media-exemplo}
\end{equation}

\guiapolitica{Use numeração e referência quando a equação precisar ser identificada no texto. A suíte matemática verifica a estabilidade tipográfica e a integração do conteúdo matemático com os motores suportados.}

\subsection{Citação longa, alíneas e enumerações}

Uma citação direta longa também é exercitada no documento de referência. O conteúdo abaixo é sintético e existe apenas para permitir a verificação geométrica do bloco:

\Enquote{A composição tipográfica de um documento científico deve ser tratada como parte de um processo reprodutível de edição, revisão e validação, mantendo conteúdo e apresentação suficientemente separados para que mudanças editoriais possam ser verificadas de forma sistemática.}
\guianormativa{ABNT NBR 10520:2023}{O bloco acima exercita a apresentação de citação direta longa descrita na Seção~\ref{sec:citacoes-referencias}: recuo de 4 cm, fonte reduzida, espaço simples e ausência de aspas no PDF final \parencite{abnt10520}. Por ser conteúdo sintético de teste, ele não recebe página ou localizador inventado; citações diretas reais devem informar o localizador aplicável quando a fonte o fornecer.}

Exemplo de alíneas e subalíneas:

\begin{ufclettereditems}
  \item primeiro item de validação;
  \item segundo item de validação:
    \begin{ufcdashedsubitems}
      \item subitem de validação.
    \end{ufcdashedsubitems}
\end{ufclettereditems}

Exemplo de enumeração comum:

\begin{enumerate}
  \item preparar o conteúdo;
  \item compilar o documento;
  \item revisar o PDF produzido.
\end{enumerate}

\subsection{Tipos bibliográficos exercitados}

O arquivo \texttt{backmatter/references.bib} reúne livro, artigo, evento, dissertação, relatório, normas e recursos eletrônicos. Campos como \texttt{isbn}, \texttt{doi}, \texttt{url}, \texttt{urldate}, \texttt{institution}, \texttt{edition}, \texttt{volume}, \texttt{number} e \texttt{pages} são usados apenas quando pertinentes ao tipo de documento.

\nocite{lamport1994,wessberg2000,vaswani2017attention,almeida2018,rfc9110}

\index{figuras}\index{gráficos}\index{quadros}\index{tabelas}\index{código-fonte}\index{algoritmos}\index{equações}

\section{Recursos do abntexto-ufc, elementos pós-textuais e revisão final}\label{sec:recursos-finais}

A classe reúne opções de perfil e módulos editoriais para que o mesmo projeto possa atender diferentes tipos de trabalho sem duplicar a lógica de formatação. Esta seção encerra o guia explicando os principais recursos e os elementos que aparecem depois do texto principal.

\subsection{Perfis de documento}

A opção \texttt{type} seleciona o perfil acadêmico. A matriz de certificação atual exercita \texttt{undergraduate-capstone}, \texttt{specialization-capstone}, \texttt{masters-thesis}, \texttt{doctoral-thesis}, \texttt{research-project} e \texttt{anonymized-research-project}.

\guiapolitica{Perfis controlam a aplicabilidade de elementos e metadados. O objetivo é evitar que o usuário tenha de remover manualmente partes da classe para adaptar um TCC a uma dissertação, tese ou projeto.}

\guianormativa{ABNT NBR 15287:2025}{Os perfis \texttt{research-project} e \texttt{anonymized-research-project} usam a estrutura de projeto de pesquisa adotada pelo projeto. O perfil anonimizado adiciona uma política técnica de ocultação de dados pessoais para processos que exijam avaliação sem identificação; essa anonimização não é apresentada como regra geral da NBR 15287 \parencite{abnt15287}.}

\subsection{Fontes e motores de compilação}

A opção \texttt{font=times|arial} seleciona a família textual desejada. A opção \texttt{strict-font} determina se a identidade literal é obrigatória ou se o ambiente pode usar fallback portátil.

\guiainstitucional{A política tipográfica UFC adotada pelo projeto trabalha com Arial ou Times New Roman. A certificação literal é executada em Windows com pdfLaTeX e LuaLaTeX; as fontes Microsoft não são redistribuídas pelo repositório.}

\guiapolitica{O documento de referência usa modo portátil para compilar em ambientes públicos. Uma compilação com fallback não deve ser confundida com evidência de identidade literal da família solicitada.}

\subsection{PDF/A}

O cabeçalho de \texttt{main.tex} declara \texttt{pdfstandard=A-2b}. A suíte verifica o PDF gerado com veraPDF e também confere incorporação de fontes.

\guiainstitucional{O fluxo de depósito acompanhado pelo projeto exige PDF/A nos casos institucionais aplicáveis. As páginas atuais do Sistema de Bibliotecas devem ser verificadas conforme o tipo de trabalho \parencite{ufcdepositotcc2026,ufcdepositopos2026}.}

\guiapolitica{A escolha do perfil específico PDF/A-2b para o documento de referência e para a certificação é uma política técnica do projeto; ela não é apresentada como uma afirmação de que toda norma ABNT exige especificamente PDF/A-2b.}

\subsection{Referências}

As referências aparecem imediatamente após o texto quando produzidas por \texttt{\string\ufcPrintReferences}. O banco bibliográfico é declarado com \texttt{\string\ufcAddBibliographyResource} e processado por Biber no fluxo recomendado.

\guianormativa{ABNT NBR 6023:2025}{Inclua na lista final as fontes de acordo com a política bibliográfica do trabalho e mantenha os dados necessários à identificação de cada documento \parencite{abnt6023}.}

\subsection{Glossário}

O glossário é um elemento opcional para definir termos especializados, siglas ou conceitos quando essa lista realmente auxiliar o leitor. O documento de referência o mantém habilitado para demonstrar a integração com o pacote \texttt{glossaries}.

\guiapolitica{Não use o glossário como duplicação automática de toda palavra técnica. Registre apenas termos cuja definição agregue valor ao trabalho e cuja repetição no corpo prejudicaria a leitura.}

\subsection{Apêndices}

Apêndices contêm material complementar elaborado pelo próprio autor e são identificados em sequência própria.

\guianormativa{estrutura pós-textual adotada pelo projeto}{Cada apêndice inicia em página própria. No escopo auditado, seu cabeçalho é centralizado, em letras maiúsculas, negrito e tamanho 12.}

\guiaexemplo{Este documento possui apêndices com material textual, questionário, código-fonte e orientação para documento externo autoral. Eles existem para exercitar diferentes comprimentos e conteúdos.}

\subsection{Anexos}

Anexos são usados para material complementar de autoria externa ou não produzido pelo autor do trabalho.

\guianormativa{estrutura pós-textual adotada pelo projeto}{Cada anexo inicia em página própria e usa a mesma política de apresentação do cabeçalho exercitada para apêndices.}

\guiaexemplo{Se um documento externo em PDF precisar ser anexado, verifique legibilidade, direitos de uso, orientação de página e compatibilidade com o arquivo final. O exemplo deste projeto documenta a rota sem redistribuir material externo desnecessário.}

\subsection{Índice}

O índice remissivo é diferente do sumário: ele organiza termos e assuntos para facilitar a localização de ocorrências no documento.

\guianormativa{ABNT NBR 6034:2004}{O índice é elemento opcional e possui norma específica de apresentação. A classe demonstra o recurso com \texttt{imakeidx} e verifica o título e sua composição no PDF final \parencite{abnt6034}.}

\subsection{Lombada}

\guianormativa{ABNT NBR 12225:2023}{A lombada possui aplicação condicional, normalmente relacionada à produção física do trabalho. O documento eletrônico de referência não a força em perfis nos quais não seja aplicável \parencite{abnt12225}.}

\subsection{Checklist antes da entrega}

Antes de entregar um trabalho real produzido a partir deste modelo:

\begin{enumerate}
  \item revise todos os campos de \texttt{\string\ufcsetup};
  \item mantenha apenas os pré-textuais aplicáveis ao seu perfil;
  \item substitua todos os exemplos e conteúdos de demonstração pelo texto acadêmico real;
  \item confirme citações e referências contra as fontes consultadas;
  \item revise figuras, tabelas, quadros, códigos, algoritmos, equações, fontes e notas;
  \item verifique glossário, apêndices, anexos e índice somente quando pertinentes;
  \item compile o documento completo com todos os processadores auxiliares necessários;
  \item examine visualmente o PDF final, inclusive quebras de página e elementos externos;
  \item execute os gates de validação disponíveis no projeto;
  \item confira as regras atuais do curso, programa e fluxo de depósito da UFC.
\end{enumerate}

\guiapolitica{Um CI verde demonstra que o corpus testado respeitou os contratos executáveis do projeto. Ele não transforma automaticamente toda regra externa em prova formal de conformidade e não substitui a revisão acadêmica e institucional do conteúdo.}

Este modelo comentado deve ser entendido como ponto de partida auditável: o usuário escreve o conteúdo acadêmico, enquanto a classe concentra a maior parte das decisões editoriais e o repositório registra a evidência usada para mantê-las.

\index{PDF/A}\index{glossário}\index{apêndice}\index{anexo}\index{índice}\index{projeto de pesquisa}


\ufcPrintReferences
\ufcPrintGlossary

\appendix{Material complementar elaborado pelo autor}
\input{backmatter/appendices/appendix-a}

\appendix{Questionário de exemplo}
\input{backmatter/appendices/appendix-b}

\appendix{Código-fonte de exemplo}
\input{backmatter/appendices/appendix-c}

\appendix{Orientação para documento externo autoral}
\input{backmatter/appendices/appendix-d}

\annex{Documento complementar externo}
\label{an:exemplo-a}

Este anexo exemplifica material não elaborado pelo autor e utilizado como documentação complementar ao trabalho.

Substitua este conteúdo pelo documento externo pertinente, preservando a identificação do anexo no comando \texttt{\string\annex} usado pelo documento principal.

\noindent\textbf{Fonte:} Instituição ou autor responsável pelo documento externo (ano). Substitua esta indicação pela fonte real do material anexado.


\annex{Orientação para documento externo em PDF}
\label{an:documento-pdf}

Para anexar um documento externo em PDF, adicione primeiro o arquivo ao projeto e então use \texttt{\string\includepdf} com um caminho existente.

A distribuição atual não inclui chamadas para PDFs ausentes, de modo que os exemplos possam ser compilados sem referências quebradas. Ao ativar um documento externo, registre também sua fonte de forma explícita no anexo e mantenha a referência bibliográfica correspondente quando aplicável.


\ufcPrintIndex

\end{document}
